%% Chapter 5 — Section 5.1: Motivating Example
%% A First Course in Monte Carlo Methods
%% D. Sanz-Alonso & O. Al-Ghattas

\section{Motivating Example}
\label{sec:5.1}

Consider a dataset $y = (y^{(1)}, \ldots, y^{(K)})$ of $K$ observations that we model as i.i.d.\
draws from a Gaussian distribution with mean $\mu$ and precision $\tau = \frac{1}{\sigma^2}$. We
take a Bayesian approach and choose the following priors for the precision $\tau = \frac{1}{\sigma^2}$
and mean $\mu$:
\begin{align*}
  \tau &\sim \GammaDist(\alpha, \beta),\\
  \mu  &\sim \Normal(m, s^2).
\end{align*}
The posterior for $\mu, \tau$ is
\[
  f(\mu, \tau \mid y)
  \propto
  \exp\!\Bigl(-\tfrac{1}{2s^2}(\mu - m)^2\Bigr)
  \exp\!\Bigl(-\tfrac{\tau}{2}\sum_{k=1}^{K}(y^{(k)} - \mu)^2\Bigr)
  \exp(-\beta\tau)\,\tau^{\alpha - 1 + \frac{K}{2}}.
\]
There is no closed form expression for the normalizing constant. Note that:
\begin{enumerate}
  \item Conditional on $\tau$,
    \[
      \mu|\tau \sim \Normal\!\left(\frac{K\bar{y}\tau + ms^{-2}}{K\tau + s^{-2}},\;
        \frac{1}{K\tau + s^{-2}}\right),
    \]
    where $\bar{y}$ is the sample average.

  \item Conditional on $\mu$,
    \[
      \tau|\mu \sim \GammaDist\!\left(\alpha + \frac{K}{2},\;
        \beta + \frac{1}{2}\sum_{k=1}^{K}(y^{(k)} - \mu)^2\right).
    \]
\end{enumerate}

Since both of these conditional distributions are standard, it seems natural to use them as
proposal distributions. Suppose that $(\mu^{(n)}, \tau^{(n)})$ are given. Inspired by the
Metropolis Hastings algorithm studied in Chapter~4, a natural idea would be to obtain
$(\mu^{(n+1)}, \tau^{(n+1)})$ as follows:
\begin{enumerate}
  \item Propose
    \[
      \mu^* \sim \Normal\!\left(\frac{K\bar{y}\tau^{(n)} + ms^{-2}}{K\tau^{(n)} + s^{-2}},\;
        \frac{1}{K\tau^{(n)} + s^{-2}}\right).
    \]

  \item Accept $\mu^{(n+1)} = \mu^*$ with probability
    $a\bigl((\tau^{(n)}, \mu^{(n)}),(\tau^{(n)}, \mu^*)\bigr)$.
    Otherwise, set $\mu^{(n+1)} = \mu^{(n)}$.

  \item Propose
    \[
      \tau^* \sim \GammaDist\!\left(\alpha + \frac{K}{2},\;
        \beta + \frac{1}{2}\sum_{k=1}^{K}(y^{(k)} - \mu^{(n+1)})^2\right).
    \]

  \item Accept $\tau^{(n+1)} = \tau^*$ with probability
    $a\bigl((\tau^{(n)}, \mu^{(n+1)}),(\tau^*, \mu^{(n+1)})\bigr)$.
    Otherwise set $\tau^{(n+1)} = \tau^{(n)}$.
\end{enumerate}

The fact that the normalizing constant of the posterior is unknown does not obstruct the
implementation of this natural scheme. Moreover, the theory we will develop in the next section
will show that the acceptance probability in steps~2 and~4 is always~1. Therefore, we could
simplify the algorithm as follows:
\begin{enumerate}
  \item Sample
    \[
      \mu^{(n+1)} \sim \Normal\!\left(\frac{K\bar{y}\tau^{(n)} + ms^{-2}}{K\tau^{(n)} + s^{-2}},\;
        \frac{1}{K\tau^{(n)} + s^{-2}}\right).
    \]

  \item Sample
    \[
      \tau^{(n+1)} \sim \GammaDist\!\left(\alpha + \frac{K}{2},\;
        \beta + \frac{1}{2}\sum_{k=1}^{K}(y^{(k)} - \mu^{(n+1)})^2\right).
    \]
\end{enumerate}

We have derived a Gibbs sampler for the posterior distribution $f(\mu, \tau \mid y)$.
